\section{Original convexity gives the same prescribed central nef class}

Fix an infinite perfect field $k$, the original projectively normal
projective subvariety $X$, a positive rational $a$, and the original
locally convex function $\varphi$ in class $a$.  The original
closed-central-curve degree theorem chooses the positive $r,d$, with
$d=ra$, the original integer $B$, the original ideals $I,J$, and their
finite-log data before any reference coordinate or curve is chosen.
Write
\[
 \rho:\mathcal N\longrightarrow\operatorname{Spec}k[t],\qquad
 z:\operatorname{Spec}k\longrightarrow\operatorname{Spec}k[t]
\]
for the same normalized simultaneous model and the actual evaluation-at-zero
map.  Retain the actual categorical scheme-theoretic fibre
\[
 Z=\mathcal N\times_{\operatorname{Spec}k[t]}\operatorname{Spec}k,
 \qquad i_0:Z\longrightarrow\mathcal N,\quad p_0:Z\longrightarrow\operatorname{Spec}k.
\]
In particular $Z$ is not replaced by its reduction or by one component.

\begin{lemma}[The actual central projection is native projective]
\label{lem:intrinsic-original-normalized-central-native-projective}
\lean{EnvelopeProperty.NumericalOriginalCentralNef.normalizedCentralFiberProjection_isProjective}
\leanok
\uses{thm:pl-normalized-polynomial-projectivity,thm:native-projectivity-from-polynomial-factorization}
For every original pair of ideals with the actual finite-log data, the
same central projection $p_0$ is projective in the native finite
quasicoherent-projective-bundle sense.  No integrality, reducedness,
normality or dimension condition on $Z$ is required.
\end{lemma}
\begin{proof}
The original normalized-model projectivity theorem supplies an actual
finite polynomial closed factorization of $\rho$.  Equip $k$ with the
actual $k[t]$-algebra structure induced by evaluation at zero.  The
existing coefficient-base-change construction gives a finite polynomial
closed factorization of precisely $p_0$: its second projection is that
of the actual pullback defining $Z$, with no changed base map.  Apply
the proved equivalence from finite polynomial factorization to native
projectivity over the same coefficient spectrum.  This argument uses
the original scheme structure of the fibre throughout.
\end{proof}

\begin{theorem}[The original prescribed central class belongs to the actual nef cone]
\label{thm:intrinsic-original-convex-central-nef}
\lean{EnvelopeProperty.NumericalOriginalCentralNef.exists_prescribedClass_originalCentralNef}
\leanok
\uses{thm:intrinsic-numerical-original-central-curve-degree,lem:intrinsic-original-normalized-central-native-projective,thm:native-numerical-nef-cone,thm:projective-central-fibre-actual-parameter-ideal,def:actual-normalized-pl-projective-twist-line,def:fc-original-curve-line-degree}
The same internally chosen $r,d,B,I,J$ and finite-log data retain their
whole original ratio-weight identity at every actual valued model
point.  For every original reference index $j$, let
\[
 M_j=\operatorname{normalizedSimultaneousTwistedRatioLine}(I,J,j,d).
\]
Then the class of the actual invertible line $i_0^*M_j$ belongs to
\[
 \operatorname{NativeRealNumericalClasses.nefCone}(p_0),
\]
the existing true real numerical quotient tested by every original
closed integral curve of $Z$.  The choices precede the universal
quantifiers over $j$ and over all those curves.  No degree-sign,
numerical-matching or nef datum is supplied as an input.
\end{theorem}
\begin{proof}
Apply the proved original-convexity closed-curve degree theorem once.
It constructs the same positive $r,d$, the relation $d=ra$, the original
$B,I,J$, their finite-log data, and the whole original ratio-weight
identity, independently of any later coordinate or curve.  Retain all
these objects unchanged.  The preceding lemma gives native projectivity
of the actual $p_0$.  The original twisted-ratio line is invertible;
ordinary pullback gives invertibility of the same $i_0^*M_j$.

Fix an actual closed integral curve $C$ of $Z$.  Its own structure map
$q:C\to\operatorname{Spec}k$ is the literal composite $C\hookrightarrow
Z\xrightarrow{p_0}\operatorname{Spec}k$; native projectivity of $p_0$
internally gives the required finite polynomial projective factorization
of $q$.  Evaluation at zero is surjective, so $z$ and its actual base
change $i_0$ are closed immersions.  Thus the original composite
\[
 c:C\hookrightarrow Z\hookrightarrow\mathcal N
\]
is a closed immersion.  The actual pullback condition, composed with
$C\hookrightarrow Z$, is precisely the original central triangle
$c\rho=qz$.  The curve's integral structure and whole dimension one
are exactly those of the original closed-curve object.

Apply the previously selected original degree inequality to this same
$C,q,c$.  It gives $\deg_q(c^*M_j)\ge0$.  Native pullback composition
gives a whole original line isomorphism
\[
 (C\hookrightarrow Z)^*(i_0^*M_j)\simeq c^*M_j.
\]
Whole Euler degree is invariant under this isomorphism, so the same
degree inequality holds for the actual curve pullback of $i_0^*M_j$.
Cast this true integer inequality to the reals.  The already proved
line-class membership criterion for the existing numerical nef cone
is exactly this collection of inequalities for all the original closed
integral curves.  It proves the asserted class membership, without
changing the line, the curve maps or the numerical carrier.  The whole
ratio-weight identity was retained verbatim from the original producer.
\end{proof}

This statement uses the existing closed-original-curve numerical
quotient.  It does not identify it with the separate numerical quotient
testing arbitrary curve maps.  It does not assert relative ampleness,
Seshadri sufficiency, uniform generation, decreasing FS approximation,
or the whole LocalPositivity or Main theorem.
